\documentclass[11pt]{article}
\usepackage[a4paper,margin=28mm]{geometry}
\usepackage{amsmath,amssymb,amsthm,mathtools,hyperref}
\newcommand{\ddc}{dd^c}
\newcommand{\PSH}{\operatorname{PSH}}
\newcommand{\MA}{\operatorname{MA}}
\newcommand{\ord}{\operatorname{ord}}
\newcommand{\PP}{\mathbb P}
\newcommand{\I}{\mathcal I}
\newcommand{\cO}{\mathcal O}
\newtheorem{theorem}{Theorem}
\newtheorem{lemma}[theorem]{Lemma}
\newtheorem{proposition}[theorem]{Proposition}
\theoremstyle{remark}\newtheorem{remark}[theorem]{Remark}
\title{Exact intrinsic traces on scalar Hopf manifolds\newline
and unbounded mass in every dimension at least three}
\author{}
\date{}
\begin{document}
\maketitle
\noindent\textbf{Working draft.} Mathematical review, independent
reconstruction and final comparison are pending. The unrestricted
non-K\"ahler trace problem remains open.
\begin{abstract}
For every scalar Hopf manifold of dimension at least three, we extend
the original smooth Xia--Zhang threefold construction to obtain
unbounded Monge--Amp\`ere mass. An exact comparison of angularly
integrated densities supplies the extension. On these targets, ambient
potentials pulled back from projective space nevertheless have exact
intrinsic traces on every smooth compact modification. The full trace
envelope descends to a closed trace on the projective base. A finite
divisorial extraction argument works in arbitrary base dimension:
blow up each divisor image and resolve the target without changing its
smooth locus, while the order of a pulled-back top form decreases.
We prove ideal-type compatibility, addition and analytic approximation,
and compute positive finite individual masses. These conclusions treat
the specified pullback data; unrestricted non-K\"ahler exactness and
ordinary-type independence are not claimed.
\end{abstract}

We use $d^c=\frac{i}{4\pi}(\bar\partial-\partial)$ and
$\ddc=\frac{i}{2\pi}\partial\bar\partial$.
Potentials have their genuine canonical representatives.
The ideal $\I(tu)$ is defined by local integrability of
$|g|^2e^{-tu}$, with no additional factor two.
An ideal restricted to a submanifold always means its coherent image
ideal. Non-pluripolar products with nonclosed smooth backgrounds use
bounded local products with smooth corrections and truncation
\cite[Section 2.1]{BGL}.

\section{Geometry and statement}
Fix an integer $n\ge3$ and put $d=n-1$. Let
\[
 M=(\mathbb C^{n+1}\setminus\{0\})/\langle z\mapsto ez\rangle,
 \qquad
 Y=H_n=(\mathbb C^n\setminus\{0\})/\langle z\mapsto ez\rangle.
\]
The inclusion of the first $n$ coordinates identifies $Y$ with a
linear submanifold of $M$. Projectivization gives the commuting square
\[
 \begin{matrix}Y&\longrightarrow&M\\
 q\downarrow&&\downarrow p\\
 B=\PP^d&\longrightarrow&\PP^n.
 \end{matrix}
\]
Both vertical maps are holomorphic submersions with connected elliptic
fiber $E=\mathbb C/(\mathbb Z+2\pi i\mathbb Z)$.
Normalize the Fubini--Study form $\kappa$ on $B$ by
$\int_B\kappa^d=1$. On $Y$, put
\[
 r=\log\|z\|^2\pmod 2,
 \quad \alpha=q^*\kappa=\ddc r,
 \quad \beta=dr\wedge d^cr,
 \quad \Omega_{a,b}=a\alpha+b\beta \quad(a,b>0).
\]
The differential of $r$ is global. The form $\beta$ has rank one,
is positive along the fibers, and $\int_E\beta=2$.
Thus $\Omega_{a,b}$ is Hermitian, $\int_Y\Omega_{a,b}^n=2nba^d$,
and $d\Omega_{a,b}=-b\,dr\wedge\alpha\ne0$.
The target is non-K\"ahler, since its first Betti number is one.
Section~\ref{sec:bmp} proves that it fails BMP for every $n\ge3$,
using the original threefold family of \cite{XZ}.

Let $U$ be genuine quasi-psh on $\PP^n$ with $\nu(U,B)=0$,
and let $u=p^*U$. Set
\[
 \mathfrak b_t=\I(tU)\cO_B,
 \qquad \mathfrak a_t=\I(tu)\cO_Y.
\]
These image ideals are nonzero. For a prime divisor $F$ over $B$,
write
\[
 \tau_B(F)=\lim_{t\to\infty}t^{-1}\ord_F\mathfrak b_t;
\]
define $\tau_Y$ from $\mathfrak a_t$ in the same way.
For any sufficiently large background $\Theta$, let
$T_\Theta$ denote the usc supremum of all $\Theta$-psh potentials
$v\le0$ satisfying $\nu(v,F)\ge\tau(F)$ on every smooth compact
modification. The definition includes all competitors.

\begin{theorem}\label{thm:main}
There is $a_0>0$, depending on the ambient curvature bound for $U$,
such that for $a\ge a_0$ and every $b>0$,
\[
 T_{\Omega_{a,b}}(u;Y)=q^*T_{a\kappa}(U;B).
 \tag{1}
\]
For every smooth compact modification $\pi:Z\to Y$ and every prime
divisor $D$ on $Z$, simultaneously,
\[
 \nu(T_{\Omega_{a,b}}(u;Y)\circ\pi,D)=\tau_Y(D).
 \tag{2}
\]
Equality (2) also holds for the trace in any Hermitian background
$\Omega\ge\Omega_{a,b}$. In the backgrounds in (1), the trace has
positive finite individual non-pluripolar mass and
\[
 \int_Y\MA_{\Omega_{a,b}}(T_{\Omega_{a,b}})
   =2nb\int_B\MA_{a\kappa}(T_{a\kappa}),
 \qquad
 2nb(a/2)^d\le\int_Y\MA_{\Omega_{a,b}}(T_{\Omega_{a,b}})
   \le2nba^d.
 \tag{3}
\]
The ideal class of the traces in (1) is independent of $a,b$ within
this range, is additive and positively homogeneous in the ambient
ideal trace data, and admits the dyadic analytic approximation
described below.
\end{theorem}

\section{Unbounded mass in all Hopf dimensions}\label{sec:bmp}
\begin{proposition}\label{prop:bmp}
For every $n\ge3$, $H_n$ fails the bounded mass property.
There are globally smooth $V_\eta$ with
$\alpha+\beta+\ddc V_\eta>0$ and total mass tending to infinity.
\end{proposition}
\begin{proof}
Use the original Xia--Zhang family, with its original heat scale
$\varepsilon_\eta=\eta^{1/\eta^2}$, not a new transport family.
Write $\xi\in\mathbb C^d$, $x=\log|\xi|^2$,
$w=\log z_0\pmod{\mathbb Z+2\pi i\mathbb Z}$ and
$\mu=e^x/(1+e^x)$ on the projective affine chart. Then
$r=w+\bar w+\log(1+e^x)$.
There are $n-2$ angular directions orthogonal to the radial direction.
Differentiating in these directions gives the identical angular
entries $A=(\mu+V_x)/\mu$. The radial/fiber block, in the frame
of \cite[Lemma 2.2]{XZ}, is
\[
 K_V=\begin{pmatrix}
 (\mu+V_{xx})/[\mu(1-\mu)]&
 (\mu+V_{x\bar w})/\sqrt{\mu(1-\mu)}\\
 (\mu+V_{wx})/\sqrt{\mu(1-\mu)}&1+V_{w\bar w}
 \end{pmatrix}.
\]
The background block has determinant one. The differential-form
expansion contains $-(1-\mu)V_x$, but its radial contribution
cancels that of the angular prefactor, giving exactly this block.
Thus the positivity conditions in every dimension are
\[
 1+V_{w\bar w}>0,\qquad \mu+V_x>0,\qquad
 Q[V]=V_{xx}+\mu-\frac{|\mu+V_{x\bar w}|^2}{1+V_{w\bar w}}>0.
\]
The source operator is $i\partial\bar\partial$; dividing both
source background forms by $2\pi$ gives ours with the same
numerical potential and these same ratios.

Normalize angular contact measure by mass one and write $dA_E$
for Haar probability measure on the elliptic fiber. The radial
Fubini--Study density on $\PP^d$ is
$d\mu^{d-1}\mu'\,dx=d\mu^d(1-\mu)\,dx$.
Since the fiber mass is two, the angularly integrated background
volume density is $2nd\mu^d(1-\mu)\,dx\,dA_E$.
Let $J_n(V;x,w)$ denote the angularly integrated smooth
Monge--Amp\`ere density. Multiplying by $A^{n-2}\det K_V$ gives
\[
 J_n(V;x,w)=\frac{n(n-1)}6(\mu+V_x)^{n-3}J_3(V;x,w).
 \tag{7}
\]
This compares integrated densities; the angular spheres in the
different dimensions each have their own probability normalization.

The original family is $V_\eta=W_\eta+g_\eta$.
It has all three strict positivity inequalities, is independent
of $x$ at the negative end, and is one constant at the positive end
\cite[Proposition 4.1]{XZ}. These endpoint properties give the
same global smooth extension with $\xi\in\mathbb C^d$.
Its construction also gives
\[
 \begin{gathered}
 \sup[W_{\eta,x}]_-=o(1),\qquad
 -m_\eta\le g_\eta'\le0,\qquad m_\eta=o(1),\\
 g_\eta''=r_\eta-m_\eta\psi,\qquad
 \int r_\eta=m_\eta,\quad\int\psi=1.
 \end{gathered}
\]
Here $r_\eta,\psi\ge0$ and $g_\eta=0$ on $x\le-2$.
Consequently $\mu+V_{\eta,x}\ge1/4$ for $x\ge0$ and small $\eta$.

The source proves $\int J_3(V_\eta)\to\infty$. To localize this
divergence, on $x<0$ one has
$W_\eta=F_\eta(\varepsilon_\eta,w)$, with no $x$ derivatives.
Since $\int_EF_{w\bar w}\,dA_E=0$, exact fiber integration gives
\[
 \int_E J_3(V_\eta;x,w)\,dA_E
   =12(\mu+g_\eta')[g_\eta''+\mu(1-\mu)].
\]
Use $0<\mu+g_\eta'\le\mu$ and
$\int|g_\eta''|\le2m_\eta$. The total mass on $x<0$ is at most
$12\int_{-\infty}^0\mu^2(1-\mu)\,dx+24m_\eta
=3/2+24m_\eta$.
The high amplitude of $F_{w\bar w}$ has been integrated exactly;
no uniform derivative bound is needed there.
It follows that $\int_{x\ge0}J_3(V_\eta)\to\infty$.
Equation (7) now gives
\[
 \int_{H_n}(\alpha+\beta+\ddc V_\eta)^n
 \ge\frac{n(n-1)}6\,4^{-(n-3)}\int_{x\ge0}J_3(V_\eta)
 \longrightarrow\infty.
\]
Every other Hermitian form $h$ dominates $c(\alpha+\beta)$
for some $c>0$. The potentials $cV_\eta$ have strictly positive
perturbed $h$, and their smooth masses dominate $c^n$ times
the preceding divergent masses. This proves BMP failure.
\end{proof}

\section{Divisorial extraction in arbitrary dimension}
\begin{lemma}\label{lem:extract}
Let $f:Z\to B$ be a dominant holomorphic map of connected smooth
compact complex manifolds, with $\dim B=d$, and let $D$ be a prime
divisor on $Z$. Either $f(D)=B$, or there is a smooth compact
modification $\sigma:B'\to B$ such that $f$ lifts holomorphically
near generic $D$ and $D$ dominates a prime divisor $F$ on $B'$.
For an integer $e\ge1$, every nonzero coherent ideal $\mathfrak b$
and every genuine quasi-psh $w$ on $B$ satisfy
\[
 \ord_D f^*\mathfrak b=e\ord_F\sigma^*\mathfrak b,
 \qquad \nu(w\circ f,D)=e\nu(w\circ\sigma,F).
 \tag{8}
\]
In the horizontal case both orders and generic Lelong numbers are zero.
\end{lemma}
\begin{proof}
Maintain a compact analytic graph proper over $Z$ which is an
isomorphism near generic $D$, with a dominant map to a smooth
compact modified base. The graph may be singular away from this
unchanged neighborhood. The divisor image $S$ is irreducible
analytic by the proper mapping theorem
\cite[Chapter II, Theorem 8.8]{Dem}.
Near generic $D$, define $J_D(f)$ as the minimum order of the
coefficients of $f^*(dy_1\wedge\cdots\wedge dy_d)$ in a smooth
source frame. Dominance gives a finite nonnegative integer,
independent of the local smooth coordinate choices.

If $S$ has codimension $c\ge2$, blow up its reduced coherent ideal.
The blowup can be singular globally, but it is smooth over the
smooth part of $S$. There use $S=\{y_1=\cdots=y_c=0\}$ and let
$m=\min_{i\le c}\ord_D f^*y_i\ge1$.
Choose $y_1$ of least order and shrink at generic $D$ so its
leading coefficient is a unit. All ratios $f^*y_i/f^*y_1$ are
holomorphic there, giving a lift to the smooth blowup chart
\[
 (y_1,\ldots,y_c)=(s,st_2,\ldots,st_c).
\]
Take the main graph of this lift. Resolve the target blowup using
a proper resolution which is an isomorphism on its whole smooth
locus \cite[Theorem 2.0.1]{Wlo}. The new target is smooth and
compact. The lift at generic $D$ is preserved, since its image
lies in that smooth locus. Taking the main graph of the resulting
meromorphic lift again preserves the generic source neighborhood.

The smooth-center Jacobian identity is
\[
 dy_1\wedge\cdots\wedge dy_d
   =s^{c-1}ds\wedge dt_2\wedge\cdots\wedge dt_c
                  \wedge dy_{c+1}\wedge\cdots\wedge dy_d.
\]
The target resolution is locally an isomorphism here. Therefore
\[
 J_D(f_{\rm new})=J_D(f)-(c-1)m\ge0.
\]
Each step strictly decreases $J$. After finitely many steps the
image has codimension one; it cannot dominate the target if its
original image was proper. This proves extraction without a
resolution of the source graphs that might change generic $D$.
Generic ideal generators along $F$ are powers of its equation
times a unit ideal, giving the first equality in (8).

At generic smooth points, tangential dominance $D\to F$ gives
rank $d-1$. Absorb a unit into a local normal coordinate using
an $e$-th root, and use the tangential base coordinates as source
coordinates. The exact local form is
\[
 (z_1,z_2,\ldots,z_d,z')\longmapsto(z_1^e,z_2,\ldots,z_d).
\]
After a smooth local correction, divisorial Siu removal writes
$w\circ\sigma=c\log|y_1|^2+R$, with $R$ psh and generic
Lelong number zero on $F$
\cite[Chapter III, (8.15)--(8.16)]{Dem}.
At a very general point of $F$ its full point Lelong number is zero.
Skoda's lemma gives integrability of $e^{-tR}$ for each $t>0$
\cite[Lemma 5.6(a)]{DemAM}.
For $e>1$ choose $1<p<e/(e-1)$. The ramification density
$|y_1|^{2/e-2}$ lies in $L^p$, so change of variables and
H\"older, using Skoda for $tp'R$, give
\[
 \int e^{-tR(z_1^e,z_2,\ldots,z_d)}dV_z
 =C_e\int e^{-tR(y)}|y_1|^{2/e-2}dV_y<\infty.
\]
Free variables contribute a harmless product factor. For $e=1$
use Fubini. A positive generic divisor coefficient $c'$ of the
residual pullback would force nonintegrability when $tc'\ge1$.
Its coefficient is therefore zero, and the logarithm gives $ec$.
This proves the Lelong equality, including when $R|_F\equiv-\infty$.
In the horizontal case $D\to B$ is a submersion at very general
points. Skoda and Fubini give zero generic coefficient; every
nonzero ideal is generically a unit there.
\end{proof}

\section{Ideal pullback and the full envelope}
\begin{lemma}\label{lem:ideals}
For a holomorphic submersion $q$ and a genuine quasi-psh potential
$v$ on its base, $\I(tq^*v)=q^*\I(tv)$ for every real $t>0$.
In the geometry above, $\mathfrak a_t=q^*\mathfrak b_t$ and
$\nu(u,Y)=0$.
\end{lemma}
\begin{proof}
Use local product coordinates $(x,y)$. The inclusion of the extended
base ideal follows by Fubini. For the reverse inclusion expand an
integrable holomorphic germ as $g(x,y)=\sum_\gamma g_\gamma(x)y^\gamma$.
The fiber coefficient estimate for holomorphic functions bounds
$|g_\gamma(x)|^2$ by a constant times the fiber integral of $|g|^2$.
Each coefficient therefore lies in $\I(tv)$ at the selected base point.
Every finite Taylor truncation belongs to its extended ideal $J$,
so $g\in J+(y)^N$ for all $N$. Krull intersection in the Noetherian
local analytic ring modulo $J$ gives $g\in J$. Smooth corrections
do not affect integrability. Apply the result to $p$ and restrict
the ideal to $Y$. Local submersion coordinates also preserve the
generic Lelong number along the inverse image of $B$.
\end{proof}

For clarity, recall the supplied closed-target trace construction.
Choose $a\kappa$ at least twice an admissible Hermitian background on
$B$ and discard finitely many indices of the ambient analytic
regularization $U_m$. The normalized ordinary analytic models
\[
 w_j=P_{a\kappa}[U_{2^j}|_B]
\]
decrease to $w=T_{a\kappa}(U;B)$.
The trace majorant lies in $\PSH(B,a\kappa/2)$ and is below $w$.
Closed mass monotonicity and the smaller-background expansion give
$\int_B\MA_{a\kappa}(w)\ge(a/2)^d$, while the closed upper bound is
$a^d$. The decreasing positive-mass model theorem
\cite[Proposition 3.1.12]{SGPT} now gives, for every divisor $F$ over $B$,
\[
 \nu(w,F)=\lim_j\nu(w_j,F)
  =\lim_j2^{-j}\ord_F\mathfrak b_{2^j}=\tau_B(F).
 \tag{6}
\]
The model and positive-mass hypotheses are essential; arbitrary
decreasing quasi-psh convergence would not imply this equality.
The closed-target construction and analytic envelope properties are
the known trace background, rather than conclusions imported to $Y$.

\begin{proof}[Proof of Theorem \ref{thm:main}]
Apply Lemma \ref{lem:extract} to $q\circ\pi$ for every target
modification $\pi:Z\to Y$. Equations (8), (6) and Lemma
\ref{lem:ideals} give
$\nu(q^*w\circ\pi,D)=\tau_Y(D)$ for each divisor, so $q^*w$ is
an exact competitor.

The compact elliptic group acts along the fibers, preserving
$\alpha,\beta$ and every $\mathfrak a_t$. It consequently preserves
the full trace competitor family, including its constraints on all
modifications. Its usc envelope is invariant and descends to a
genuine quasi-psh potential on $B$ by local sections. For a basic
potential $q^*v$, the background inequality is equivalent to
$a\kappa+\ddc v\ge0$. Indeed wedge the background inequality
with the smooth semipositive $\beta$ and push forward by $q$.
Since $\beta^2=0$ and its fiber integral is two, the resulting
positive current is $2(a\kappa+\ddc v)$. Conversely add $b\beta\ge0$.

Every base modification $B'\to B$ lifts to the smooth compact
modification $Y\times_B B'\to Y$. The inverse image of a prime divisor
$F$ is prime, since the elliptic fibers are connected; at generic
points its multiplicity is one. Thus the descended trace envelope
satisfies every base constraint, and is at most $w$. The opposite
inequality follows from the exact competitor $q^*w$. This proves
(1) and (2). For $\Omega\ge\Omega_{a,b}$ the same exact competitor
lies below $T_\Omega$, so
$\tau_Y(D)\le\nu(T_\Omega,D)\le\nu(q^*w,D)=\tau_Y(D)$.

For $v_j=\max(v,-j)$ on the base, actual bounded local products give
\[
 (\Omega_{a,b}+\ddc q^*v_j)^n
   =nb\,\beta\wedge q^*(a\kappa+\ddc v_j)^d,
\]
since the horizontal current has rank at most $d$ and $\beta^2=0$.
Multiply by $\mathbf1_{\{v>-j\}}$ and pass to the non-pluripolar
limit. Integrating over the fibers gives the factor $2nb$ in (3).
The closed-base bounds above give its two inequalities.
These are individual masses, not extremal volumes of all ordinary
representatives.

Different sufficiently large base backgrounds produce potentials
with the same numbers (6); the closed K\"ahler ideal criterion
\cite[Proposition 3.2.1]{SGPT} gives the same ideals at every real
weight. Lemma \ref{lem:ideals} transfers this to (1).
Intrinsic trace data are additive and positively homogeneous.
On the base, the sum of two exact realizers has the sum of their
divisorial numbers. The same criterion identifies its ideal class
with the trace for $U_1+U_2$, in a background with sufficient
curvature room. Positive scaling is identical. Pull back to $Y$.
This uses the ideal criterion on the projective base, and does not
assume a global valuative criterion on an arbitrary non-K\"ahler target.
\end{proof}

\section{Approximation and a point trace}
For the dyadic $w_j$ above, put $v_j=q^*w_j$. Compact elliptic
invariance applied to the full ordinary obstacle envelopes shows,
for $h\in\PSH(B,a\kappa)$,
\[
 P_{\Omega_{a,b}}[q^*h]=q^*P_{a\kappa}[h].
\]
At each finite obstacle level the competitor family is invariant;
its full supremum descends, and ordinary order is equivalent on
the base by surjectivity. Thus $v_j$ are the full ordinary analytic
models of the pulled-back restricted approximants. They decrease
to (1) in $L^1(Y)$: decreasing quasi-psh convergence is genuine,
or one can integrate a fixed invariant volume along the fibers.
For every divisor $D$ over $Y$, (6) and (8) yield
\[
 \nu(v_j,D)\longrightarrow\tau_Y(D).
\]
The same potentials pulled back to any target modification converge
in $L^1$ there, by quasi-psh compactness and uniqueness off the
exceptional set. This assertion supplies no convergence theorem
for extremal non-K\"ahler singularity volumes or capacities.

A concrete trace involving exceptional divisors is obtained from
\[
 U([z_0:\cdots:z_n])
  =\log\frac{\sum_{i=0}^{n-2}|z_i|^2}{\sum_{i=0}^{n}|z_i|^2},
 \qquad B=\{z_n=0\}.
\]
Its generic number along $B$ is zero, and its restricted image
ideals are
\[
 \mathfrak b_t=\mathfrak m_P^{\max(\lfloor t-d+1\rfloor,0)},
 \qquad P=[0:\cdots:0:1]\in B.
\]
The formula follows from the normal radial integrability condition
$k>t-d$ in $d$ normal coordinates to the ambient singular line.
Consequently $\tau_Y(D)=\ord_Dq^*\mathfrak m_P$ on every model.
The genuine potential $q^*(U|_B)$ is an explicit exact realizer,
singular along the elliptic fiber over $P$. For $a>1$,
\[
 \int_Y\MA_{\Omega_{a,b}}(q^*(U|_B))=2nb(a^d-1).
\]
On the point blowup the horizontal current, after removing the
exceptional divisor $F$, has class $aH-F$ and smooth semipositive
representative $(a-1)\sigma^*\kappa+\eta$. Here $\eta$ is the
pullback of the Fubini--Study form by the resolved projection to
$\PP^{d-1}$. The mixed intersections of $H$ and $F$ vanish,
$H^d=1$ and $F^d=(-1)^{d-1}$, giving $(aH-F)^d=a^d-1$.
Its smooth product off $F$ is the pullback of the actual
non-pluripolar product, and a smooth measure charges no divisor.
The displayed potential need not equal the normalized trace envelope.

\begin{remark}
The theorem treats arbitrary genuine ambient $U$ with the stated
generic condition, including nonanalytic data and cases where
$U|_B$ is identically $-\infty$; image ideals, not naive restriction,
define the trace. The base has arbitrary dimension $d\ge2$. No finite selection
of exceptional divisors replaces the simultaneous constraint.
Ideal-type compatibility is asserted for the invariant backgrounds
in (1). Exactness in larger arbitrary backgrounds follows by the
sandwich argument, but ordinary-type independence is not claimed.
\end{remark}

\begin{thebibliography}{9}
\bibitem{BGL} S. Boucksom, V. Guedj and C. H. Lu,
\emph{Volumes of Bott--Chern classes}, arXiv:2406.01090.
\bibitem{Dem} J.-P. Demailly, \emph{Complex Analytic and Differential
Geometry}, Harvard-hosted June 2012 manuscript.
\bibitem{DemAM} J.-P. Demailly, \emph{Analytic Methods in Algebraic
Geometry}, Surveys of Modern Mathematics 1, 2012.
\bibitem{SGPT} M. Xia, \emph{Singularities in global pluripotential
theory}, lecture manuscript, October 4, 2026 version,
\url{https://mingchenxia.github.io/Lectures/SGPT_final.pdf}.
\bibitem{Wlo} J. W\l{}odarczyk,
\emph{Resolution of singularities of analytic spaces},
Proceedings of the 15th G\"okova Geometry--Topology Conference
(2008), 31--63, published 2009,
\url{https://gokovagt.org/proceedings/2008/ggt08-wlodarczyk.pdf}.
\bibitem{XZ} M. Xia and K. Zhang,
\emph{A counterexample to the bounded mass property},
arXiv:2608.21053v1.
\end{thebibliography}
\end{document}
